\documentclass[11pt,a4paper]{article}

\usepackage[utf8]{inputenc}
\usepackage[T1]{fontenc}
\usepackage[french]{babel}
\usepackage{geometry}
\usepackage{xcolor}
\usepackage{longtable}
\usepackage{array}
\usepackage{hyperref}
\usepackage{fancyvrb}

\geometry{margin=2.2cm}
\hypersetup{
  colorlinks=true,
  linkcolor=blue,
  urlcolor=blue
}

\title{Documentation technique -- Base GLPI Multi-Sites}
\author{Projet TAD -- CY Tech}
\date{}

\begin{document}

\maketitle
\tableofcontents
\newpage

\section{Vue d'ensemble}

Le projet modelise une base GLPI simplifiee pour deux sites : \texttt{CERGY} et \texttt{PAU}. Le modele est centre sur la table \texttt{assets}, qui regroupe les ordinateurs, imprimantes, telephones, peripheriques, moniteurs et equipements reseau.

La BDDR est simulee sur une seule instance Oracle avec deux schemas locaux :

\begin{itemize}
  \item \texttt{GLPI\_CERGY}
  \item \texttt{GLPI\_PAU}
\end{itemize}

Ces schemas exposent des vues filtrees par \texttt{sites.site\_code}, puis des DB links locaux permettent de montrer une interrogation inter-sites.

\section{Diagrammes}

\begin{longtable}{|p{5cm}|p{8cm}|}
\hline
\textbf{Diagramme} & \textbf{Fichier} \\
\hline
Architecture & \texttt{docs/diagrams/architecture.svg} \\
\hline
MCD & \texttt{docs/diagrams/mcd.svg} \\
\hline
UML & \texttt{docs/diagrams/uml.svg} \\
\hline
MLD textuel & \texttt{docs/mld.md} \\
\hline
\end{longtable}

\section{Tables}

\subsection{\texttt{sites}}

Represente les sites ou structures organisationnelles. Les sites physiques sont identifies par \texttt{site\_code = 'CERGY'} ou \texttt{site\_code = 'PAU'}.

\begin{longtable}{|p{4cm}|p{10cm}|}
\hline
\textbf{Champ} & \textbf{Role} \\
\hline
\texttt{id} & Identifiant du site. \\
\hline
\texttt{name} & Nom court. \\
\hline
\texttt{parent\_site\_id} & Site parent, pour une hierarchie simple. \\
\hline
\texttt{full\_name} & Nom complet lisible. \\
\hline
\texttt{site\_code} & Code de distribution des donnees. \\
\hline
\texttt{created\_at}, \texttt{updated\_at} & Dates techniques. \\
\hline
\end{longtable}

\subsection{\texttt{locations}}

Emplacements physiques des utilisateurs et materiels : batiments, salles, bureaux.

\subsection{\texttt{manufacturers}}

Referentiel des fabricants : Dell, HP, Lenovo, Cisco, Epson, Apple, Samsung, Ubiquiti.

\subsection{\texttt{states}}

Referentiel d'etat des materiels : en service, stock, maintenance, hors service, a renouveler.

\subsection{\texttt{users}}

Utilisateurs metier, demandeurs, techniciens et responsables.

\begin{longtable}{|p{4cm}|p{10cm}|}
\hline
\textbf{Champ} & \textbf{Role} \\
\hline
\texttt{login} & Identifiant unique. \\
\hline
\texttt{site\_id} & Site principal. \\
\hline
\texttt{location\_id} & Localisation principale. \\
\hline
\texttt{supervisor\_user\_id} & Responsable hierarchique. \\
\hline
\texttt{is\_active} & Activation du compte. \\
\hline
\end{longtable}

\subsection{\texttt{profiles}}

Profils applicatifs simplifies : administrateur, technicien, manager site, utilisateur.

\subsection{\texttt{profiles\_users}}

Association entre utilisateur, profil et site. Le champ \texttt{is\_recursive} permet d'etendre le profil a des sous-sites.

\subsection{\texttt{groups}}

Groupes fonctionnels par site : support IT, equipe reseau, administration.

\subsection{\texttt{groups\_users}}

Association entre utilisateurs et groupes. Le champ \texttt{is\_manager} indique un responsable de groupe.

\subsection{\texttt{assets}}

Table centrale de l'inventaire. Elle evite de multiplier les tables par type de materiel.

\begin{longtable}{|p{4cm}|p{10cm}|}
\hline
\textbf{Champ} & \textbf{Role} \\
\hline
\texttt{asset\_type} & Type fonctionnel : \texttt{COMPUTER}, \texttt{PRINTER}, etc. \\
\hline
\texttt{serial\_number} & Numero de serie, unique par site. \\
\hline
\texttt{owner\_user\_id} & Utilisateur principal. \\
\hline
\texttt{technician\_user\_id} & Technicien responsable. \\
\hline
\texttt{manufacturer\_id} & Fabricant. \\
\hline
\texttt{state\_id} & Etat du materiel. \\
\hline
\end{longtable}

\subsection{\texttt{ticket\_categories}}

Referentiel des categories support : incident materiel, demande reseau, installation, acces utilisateur, gestion du stock.

\subsection{\texttt{tickets}}

Demandes et incidents lies a un materiel.

\begin{longtable}{|p{4cm}|p{10cm}|}
\hline
\textbf{Champ} & \textbf{Role} \\
\hline
\texttt{asset\_id} & Materiel concerne. \\
\hline
\texttt{requester\_user\_id} & Demandeur. \\
\hline
\texttt{assigned\_group\_id} & Groupe support affecte. \\
\hline
\texttt{status} & Cycle de vie : nouveau, assigne, en cours, resolu, clos. \\
\hline
\texttt{priority} & Priorite. \\
\hline
\texttt{description}, \texttt{resolution} & Textes longs en CLOB. \\
\hline
\end{longtable}

\subsection{\texttt{ticket\_users}}

Techniciens affectes a un ticket. Un trigger verifie que l'utilisateur affecte possede le profil \texttt{Technicien}.

\subsection{\texttt{ticket\_followups}}

Historique textuel d'un ticket : demande initiale, diagnostic, resolution.

\subsection{\texttt{network\_ports}}

Ports reseau rattaches aux assets. Un equipement peut avoir plusieurs ports.

\subsection{\texttt{ip\_networks}}

Sous-reseaux IP d'un site : utilisateurs, serveurs, imprimantes, Wi-Fi. Les colonnes \texttt{vlan\_name} et \texttt{vlan\_tag} representent le VLAN sans table supplementaire.

\subsection{\texttt{ip\_addresses}}

Adresses IP affectees aux ports reseau.

\subsection{\texttt{audit\_log}}

Journal d'audit alimente par le trigger \texttt{TRG\_AUDIT\_ASSETS}.

\subsection{\texttt{archives\_materiel}}

Archive fonctionnelle des assets supprimes, alimentee avant suppression par \texttt{TRG\_ARCHIVE\_ASSET\_DELETE}.

\section{Vues metier}

\begin{longtable}{|p{5cm}|p{9cm}|}
\hline
\textbf{Vue} & \textbf{Utilite} \\
\hline
\texttt{V\_INVENTAIRE\_COMPLET} & Inventaire enrichi avec site, localisation, proprietaire, fabricant et etat. \\
\hline
\texttt{V\_MATERIEL\_PAR\_SITE} & Comptage des materiels par site et type. \\
\hline
\texttt{V\_UTILISATEURS\_PROFILS} & Liste des utilisateurs avec profils et sites d'application. \\
\hline
\texttt{V\_TOPOLOGIE\_RESEAU} & Vue des ports, IP, sous-reseaux et VLAN. \\
\hline
\texttt{V\_STATISTIQUES\_SITE} & Indicateurs par site. \\
\hline
\texttt{V\_MATERIEL\_RECENT} & Materiels modifies dans les 30 derniers jours. \\
\hline
\texttt{V\_TICKETS\_SUPPORT} & Tickets enrichis avec demandeur, techniciens, groupe, categorie et materiel. \\
\hline
\end{longtable}

\section{Fonctions PL/SQL}

\begin{longtable}{|p{5cm}|p{4cm}|p{5cm}|}
\hline
\textbf{Fonction} & \textbf{Parametres} & \textbf{Role} \\
\hline
\texttt{FN\_COMPTER\_MATERIEL\_SITE} & \texttt{p\_site\_code} & Retourne le nombre d'assets d'un site. \\
\hline
\texttt{FN\_CALCULER\_TAUX\_UTILISATION} & \texttt{p\_asset\_id} & Calcule le pourcentage de ports avec IP. \\
\hline
\texttt{FN\_OBTENIR\_SITE\_ENTITE} & \texttt{p\_site\_id} & Remonte la hierarchie pour retrouver le code site. \\
\hline
\texttt{FN\_VERIFIER\_QUOTA\_MATERIEL} & \texttt{p\_site\_code}, \texttt{p\_quota\_max} & Indique si le quota est atteint. \\
\hline
\end{longtable}

\section{Procedures PL/SQL}

\begin{longtable}{|p{5cm}|p{9cm}|}
\hline
\textbf{Procedure} & \textbf{Role} \\
\hline
\texttt{SP\_TRANSFERT\_MATERIEL} & Simule le transfert inter-sites d'un materiel. \\
\hline
\texttt{SP\_AFFECTER\_PROFIL} & Cree ou met a jour l'affectation d'un profil utilisateur. \\
\hline
\texttt{SP\_CREER\_TICKET\_MATERIEL} & Cree un ticket sur un asset. \\
\hline
\texttt{SP\_ASSIGNER\_TECH\_TICKET} & Affecte un technicien a un ticket existant. \\
\hline
\texttt{SP\_INVENTAIRE\_SITE} & Affiche un resume par site. \\
\hline
\texttt{SP\_NETTOYER\_ARCHIVES} & Supprime les archives trop anciennes. \\
\hline
\texttt{SP\_GENERER\_JEU\_DE\_TEST} & Genere les donnees de test enrichies. \\
\hline
\texttt{SP\_REPLIQUER\_REFERENTIELS} & Procedure de demonstration BDDR locale. \\
\hline
\end{longtable}

\section{Triggers}

\begin{longtable}{|p{5cm}|p{3cm}|p{6cm}|}
\hline
\textbf{Trigger} & \textbf{Table} & \textbf{Role} \\
\hline
\texttt{TRG\_AUTO\_DATE\_MOD\_ASSETS} & \texttt{assets} & Met a jour \texttt{updated\_at}. \\
\hline
\texttt{TRG\_AUTO\_DATE\_MOD\_USERS} & \texttt{users} & Met a jour \texttt{updated\_at}. \\
\hline
\texttt{TRG\_AUTO\_DATE\_MOD\_NETPORTS} & \texttt{network\_ports} & Met a jour \texttt{updated\_at}. \\
\hline
\texttt{TRG\_AUTO\_DATE\_MOD\_TICKETS} & \texttt{tickets} & Met a jour les dates du workflow. \\
\hline
\texttt{TRG\_CHECK\_TICKET\_USER\_TECH} & \texttt{ticket\_users} & Interdit d'affecter un non-technicien. \\
\hline
\texttt{TRG\_AUDIT\_ASSETS} & \texttt{assets} & Journalise les changements. \\
\hline
\texttt{TRG\_ARCHIVE\_ASSET\_DELETE} & \texttt{assets} & Archive avant suppression. \\
\hline
\texttt{TRG\_CASCADE\_SITE\_CODE} & \texttt{sites} & Propage un changement de code site. \\
\hline
\end{longtable}

\section{Clusters et index}

Les clusters physiques sont crees dans \texttt{02\_schema\_tables.sql}.

\begin{longtable}{|p{5cm}|p{3cm}|p{6cm}|}
\hline
\textbf{Cluster} & \textbf{Cle} & \textbf{Tables} \\
\hline
\texttt{cluster\_user\_rights} & \texttt{user\_id} & \texttt{users}, \texttt{profiles\_users}, \texttt{groups\_users} \\
\hline
\texttt{cluster\_site\_structure} & \texttt{site\_id} & \texttt{sites}, \texttt{locations}, \texttt{groups}, \texttt{ip\_networks} \\
\hline
\texttt{cluster\_network\_port\_ips} & \texttt{network\_port\_id} & \texttt{network\_ports}, \texttt{ip\_addresses} \\
\hline
\end{longtable}

\section{BDDR locale}

Le fichier \texttt{07\_bddr.sql} cree deux schemas de simulation :

\begin{itemize}
  \item \texttt{GLPI\_CERGY} expose les vues du fragment Cergy ;
  \item \texttt{GLPI\_PAU} expose les vues du fragment Pau.
\end{itemize}

Chaque schema peut interroger l'autre via un DB link local :

\begin{itemize}
  \item \texttt{lien\_pau} depuis \texttt{GLPI\_CERGY} ;
  \item \texttt{lien\_cergy} depuis \texttt{GLPI\_PAU}.
\end{itemize}

Les vues globales \texttt{V\_ASSETS\_GLOBAL}, \texttt{V\_USERS\_GLOBAL}, \texttt{V\_STATS\_GLOBAL} et \texttt{V\_TICKETS\_GLOBAL} font des \texttt{UNION ALL} entre le fragment local et le fragment distant.

\section{Donnees de test}

Le jeu de test genere par \texttt{09\_test\_data.sql} contient par defaut :

\begin{itemize}
  \item 2 sites ;
  \item 6 localisations ;
  \item 300 utilisateurs ;
  \item 1400 assets ;
  \item plusieurs groupes, profils, categories, fabricants et etats ;
  \item 8 sous-reseaux IP ;
  \item des ports reseau et adresses IP ;
  \item des tickets, affectations et suivis.
\end{itemize}

\end{document}
